
\documentclass[11pt]{article}
\usepackage[margin=1in]{geometry}
\usepackage{parskip}
\usepackage{hyperref}
\usepackage{enumitem}
\setlist[itemize]{topsep=2pt,itemsep=2pt,parsep=0pt,partopsep=0pt}
\setlist[enumerate]{topsep=2pt,itemsep=2pt,parsep=0pt,partopsep=0pt}
\hypersetup{colorlinks=true,linkcolor=blue,urlcolor=blue}

\title{Lecture 4 Speaker Handout\\[-2pt]\large Inflation Dynamics (Old Keynesian) + Cochrane Ch.3--4 Pointers}
\author{Hui-Jun Chen}
\date{Spring 2026}

\begin{document}
\maketitle

\section{From price level movements to inflation paths}

\subsection*{Title page}
\begin{itemize}
\item Last time: AD--AS diagnosed demand vs supply shocks in $(Y,P)$.
\item Today: turn that into a \textbf{time path} for inflation $\pi_t$.
\item Key move: replace ``one-shot price level change'' with a \textbf{law of motion} for inflation.
\end{itemize}

\subsection*{Price level $P_t$ vs inflation $\pi_t$}
\begin{itemize}
\item AD--AS tells us where the \textbf{price level} lands after a shock.
\item But inflation is the \textbf{change} in the price level: $\pi_t \equiv \Delta\log P_t$.
\item One-time jump in $P$ $\neq$ inflation staying high; persistence requires a dynamic mechanism.
\end{itemize}

\subsection*{Output gap as ``pressure''}
\begin{itemize}
\item Define $Y^{FE}$ (full-employment / natural output) from the RBC supply block.
\item Output gap $x_t \equiv y_t - y_t^{FE}$ summarizes slack vs tightness.
\item Most inflation models: inflation = expected inflation + (slack term) + shocks.
\end{itemize}

\subsection*{AD--AS recap (memory anchor figure)}
\begin{itemize}
\item Demand shock: $Y$ and $P$ move together.
\item Supply shock: $Y$ and $P$ move opposite (stagflation).
\item Now we write equations that generate the \emph{paths} behind these movements.
\end{itemize}

\section{Old-Keynesian inflation dynamics: the missing equation}

\subsection*{Adaptive expectations (benchmark)}
\begin{itemize}
\item Benchmark: expected inflation equals last period's inflation, $\pi_t^e=\pi_{t-1}$.
\item Captures inertia in a transparent way; later we upgrade to forward-looking expectations (NK).
\end{itemize}

\subsection*{Backward-looking Phillips curve}
\begin{itemize}
\item Missing equation: inflation changes with the output gap:
\[
\pi_t=\pi_{t-1}+\kappa x_t+u_t.
\]
\item If $x_t>0$ (overheating) inflation \emph{accelerates}; if $x_t<0$ inflation \emph{decelerates}.
\item $u_t$ is cost-push: inflation can rise even without a boom (bottlenecks, oil, supply disruptions).
\end{itemize}

\subsection*{IS / AD in gap form (demand side)}
\begin{itemize}
\item Reduced-form demand:
\[
x_t=-\sigma(r_t-r_t^{n})+\varepsilon_t^{d}.
\]
\item Real rate above natural rate cools demand; $\varepsilon_t^{d}$ shifts demand.
\end{itemize}

\subsection*{Policy rule + Fisher}
\begin{itemize}
\item Modern policy: interest-rate reaction rule (Taylor-style).
\item Fisher links nominal and real rates: $r_t=i_t-\pi_t^e$.
\item With adaptive expectations, $\pi_t^e=\pi_{t-1}$ is predetermined $\Rightarrow$ tightening raises real rate immediately.
\end{itemize}

\section{Algebra: derive the doctrine from equations}

\subsection*{Algebra 1: express the real-rate gap in observables}
\begin{itemize}
\item Work with inflation deviation from target $\tilde\pi_t$.
\item Combine Taylor rule + Fisher + $\pi_t^e=\pi_{t-1}$:
\[
r_t-r^n=\phi_\pi\tilde\pi_t+\phi_x x_t-\tilde\pi_{t-1}.
\]
\item Interpretation: policy response to inflation and slack maps into the real-rate gap.
\end{itemize}

\subsection*{Algebra 2: solve IS for $x_t$}
\begin{itemize}
\item Substitute the real-rate gap into IS.
\item Collect terms because the rule also depends on $x_t$:
\[
x_t=\frac{-\sigma\phi_{\pi}\tilde\pi_t+\sigma\tilde\pi_{t-1}+\varepsilon_t^{d}}{1+\sigma\phi_x}.
\]
\item Read: higher current inflation triggers tighter policy $\Rightarrow$ smaller $x_t$.
\end{itemize}

\subsection*{Algebra 3: inflation becomes AR(1) under policy feedback}
\begin{itemize}
\item Plug $x_t$ into $\tilde\pi_t=\tilde\pi_{t-1}+\kappa x_t+u_t$.
\item Rearrange to get:
\[
\tilde\pi_t
=
a\,\tilde\pi_{t-1}
+b\,\varepsilon_t^{d}
+c\,u_t,
\]
where $a$ depends on $(\kappa,\sigma,\phi_\pi,\phi_x)$.
\item Key: stability $\Leftrightarrow a<1$; policy parameters determine whether inflation mean-reverts.
\end{itemize}

\subsection*{System summary: Old-Keynesian 3-equation core}
\begin{itemize}
\item Demand: $x_t=-\sigma(r_t-r_t^{n})+\varepsilon_t^{d}$.
\item Inflation: $\pi_t=\pi_{t-1}+\kappa x_t+u_t$.
\item Policy: Taylor rule + Fisher + adaptive expectations.
\item This is ``AD--AS with dynamics'': it produces time paths for $(x_t,\pi_t)$.
\end{itemize}

\section{Implications you can state cleanly}

\subsection*{Why ``raise rates to lower inflation'' works here}
\begin{itemize}
\item Tightening $\Rightarrow$ real rate rises (because expectations are sticky).
\item IS: higher real rate lowers $x_t$ (slack).
\item Phillips: slack makes inflation decelerate: $\pi_t-\pi_{t-1}=\kappa x_t+u_t$.
\item Disinflation is gradual because inflation inherits $\pi_{t-1}$.
\end{itemize}

\subsection*{Taylor principle (derived)}
\begin{itemize}
\item In AR(1) form, stability is $a<1$.
\item $\phi_\pi>1$ tends to make $a<1$ (mean reversion); $\phi_\pi=1$ gives a unit-root style case; $\phi_\pi<1$ can be unstable.
\item Intuition: to cool inflation you need the real rate to rise when inflation rises.
\end{itemize}

\subsection*{Back-of-the-envelope: ``how much slack to disinflate?''}
\begin{itemize}
\item Rearranged Phillips curve:
\[
x_t=\frac{\pi_t-\pi_{t-1}-u_t}{\kappa}.
\]
\item Smaller $\kappa$ (flatter Phillips) implies larger output gaps for a given disinflation.
\item Cost-push ($u_t>0$) makes disinflation harder even with slack.
\end{itemize}

\subsection*{Demand vs supply inflation (connect to Lecture 3)}
\begin{itemize}
\item Demand shock $\varepsilon_t^{d}>0$: raises $x_t$, then inflation rises via $\kappa x_t$ (AD shift).
\item Cost-push $u_t>0$: inflation rises directly even if $x_t$ is small (AS shift).
\item Same taxonomy as AD--AS, now with dynamics.
\end{itemize}

\subsection*{A compact disinflation narrative}
\begin{enumerate}
\item Policy raises $i_t$.
\item With sticky expectations, $r_t$ rises.
\item Demand cools: $x_t<0$.
\item Inflation decelerates over time: $\pi_t$ drifts down.
\end{enumerate}

\section{Cochrane Ch.3--4 pointers (what to emphasize)}

\subsection*{Cochrane Ch.3 (Monetary policy) in one paragraph}
\begin{itemize}
\item Under interest-rate targeting, ``rates control inflation'' is not automatic: you must specify the broader regime.
\item In particular, whether fiscal policy provides backing matters for price-level determination.
\item The Old-Keynesian model delivers the doctrine because of backward-looking expectations + slack channel; Cochrane asks when that survives stronger expectation assumptions.
\end{itemize}

\subsection*{Cochrane Ch.4 (Distinguishing theories) in one paragraph}
\begin{itemize}
\item Different theories imply different anchors and different predictions under interest-rate pegs, Taylor rules, and ZLB episodes.
\item Old Keynesian: backward-looking inflation dynamics; policy stabilizes via real-rate channel.
\item New Keynesian: forward-looking pricing shifts the role of expectations; determinacy/equilibrium selection become central.
\item Fiscal theory: valuation of nominal government liabilities anchors the price level; fiscal regime matters.
\end{itemize}

\subsection*{Interest-rate peg thought experiment}
\begin{itemize}
\item Fix $i_t$ for a long time: what happens differs across models.
\item This is a clean way to motivate why we need to be explicit about expectations and the nominal anchor.
\end{itemize}

\section*{Next lecture (what changes)}
\begin{itemize}
\item Replace adaptive expectations with forward-looking expectations.
\item Replace backward-looking Phillips with NK Phillips curve.
\item Revisit ``Taylor principle'' and stability in that environment.
\end{itemize}

\end{document}
